\section{Definitions and theory}
\label{sec:theory}

\subsection{Deadliness of one passage}
Deadliness is the perceived intensity of the worst single inhalation, $I = k\,\mathrm{OAV}^n$, with the OAV taken as the
flow-weighted mean over one inhalation. Perceived intensity saturates within one sniff
\cite{laing1983,laing1985}, still integrates within it \cite{frasnelli2006}, and the sniff is the unit of olfactory
processing \cite{uchida2003,kepecs2006}; odour-impact regulation also works per breath \cite{schauberger2012}.
Cumulative exposure is a secondary metric.

\subsection{Linear transport (Theorem 1)}
\gap{Statement and proof.} For a passive scalar, $C = G * q$. The time-integrated dose does not depend on the release
history, and a slower release never has a higher peak (Young's inequality), for any Green's function $G$.

\subsection{Loss of release memory (Theorem 2)}
\gap{Statement.} Differences between release modes decay with time and distance; in a closed room both tend to
$M/V_\mathrm{room}$.

\subsection{Where the theorems fail}
Only release momentum and buoyancy break linearity: near-field jet pre-dilution, bounded by the Silencer number
$\mathrm{Si} = u_0/u_h$, and changes of routing, of either sign.

\subsection{Selection (Theorem 3)}
\gap{Statement and proof.} If a passage is remembered when heard ($H$) or smelled ($S$), then
$P(S\mid\neg H, N) = 1$ and $P(S\mid H, N) = P(S\mid H)$, whatever the joint law of $H$ and $S$
\cite{berkson1946,snoep2014}. An observer who notices every passage (the emitter) has no bias.
